\section{Projects for Chapter 2: Discovering Formalism}

This chapter is less about numerical simulation and more about making abstract structure visible. Two projects are especially useful: an event-algebra playground and a finite sigma-algebra generator.

\subsection*{Project: event algebra playground}

Use a small finite universe displayed as a grid of elementary outcomes. The user should define events \(A\) and \(B\) by clicking outcomes. The application should then display union, intersection, complements, differences, and symmetric difference.

\begin{examplebox}
\textbf{Prompt to give an AI coding assistant}

\begin{quote}\small
Create an event-algebra playground on a finite sample space containing between eight and twenty elementary outcomes. Let the user select events \(A\) and \(B\). Display \(A\cup B\), \(A\cap B\), \(A^c\), \(B^c\), \(A\setminus B\), and the symmetric difference using consistent highlighting.

Next to every set operation, display the corresponding logical statement using words such as and, or, not, and but not. Add buttons that verify De Morgan's laws by comparing both sides outcome by outcome. Include a panel that writes the selected event in roster notation.
\end{quote}
\end{examplebox}

\subsection*{Project: sigma-algebra generated by a partition}

A finite sigma-algebra can be visualised without advanced measure theory. Begin with a partition of a finite universe into blocks. The application should construct all unions of those blocks and explain why no other subsets belong to the generated sigma-algebra.

\begin{examplebox}
\textbf{Prompt to give an AI coding assistant}

\begin{quote}\small
Build an interactive finite sigma-algebra generator. Display a finite universe and let the user partition it into blocks. Generate the collection of all unions of complete blocks. Show that the empty set and the universe are included, that complements remain in the collection, and that unions of listed events remain in the collection.

Allow the user to propose an additional subset. Explain whether it belongs to the generated sigma-algebra and, if it does not, identify the block that it splits. Keep the explanation elementary and connect the construction with the idea that events represent distinctions the model is able to observe.
\end{quote}
\end{examplebox}

A frequency panel may be added as motivation: repeated trials can produce relative frequencies for selected events. The program should state explicitly that frequency patterns motivate probability but do not prove the axioms.

\begin{summarybox}
The most useful visualisation in this chapter is structural:
\[
\text{logical statement}
\longleftrightarrow
\text{event}
\longleftrightarrow
\text{set operation}.
\]
\end{summarybox}
